\documentclass[12pt]{article}
\usepackage{graphicx}
\usepackage{hyperref}
\usepackage{amsmath}
\usepackage{amssymb}
\usepackage{siunitx}
\usepackage[numbers, square, sort&compress]{natbib}

\title{Chapter 8: Analysing Results from Molecular Docking}
\author{Elvis Martis}
\date{}

\begin{document}

\maketitle

\section{Introduction}

Molecular docking is now a routine component of structure-based drug discovery pipelines, used at scales ranging from hypothesis-driven pose analysis of a handful of ligands to ultra–large virtual screens of hundreds of millions of compounds. Despite this apparent maturity, interpreting and validating docking results remains non-trivial and is often the weakest link in otherwise sophisticated workflows. Numerous benchmarking studies have shown that while many docking programs can often recover crystallographic binding modes with reasonable accuracy, their ability to rank ligands by affinity or to drive robust hit discovery is far more limited and system-dependent\cite{CH8_Warren2006,CH8_Kitchen2004,CH8_Cheng2009}.

This chapter focuses on practical aspects of analysing and validating docking results. We emphasise not only the good practices followed by experienced practitioners, but also the common failures and negative examples encountered in the literature and in early–stage projects. The goal is to provide a realistic view of what docking can and cannot do, and to design and interpret calculations so they educate rather than mislead.

\section{Docking Outputs}

A typical small-molecule docking calculation produces one or more putative binding modes, also called poses, for each ligand in a defined binding site. This is accompanied by a numerical score intended to correlate with binding affinity. More realistically, the score must rank ligands relative to one another based on their binding affinity\cite{CH8_Kitchen2004,CH8_Huang2011}. Many modern workflows also provide additional data, such as:

\begin{itemize}
  \item Decomposition of interaction into van der Waals, Coulombic, desolvation, ionic, and hydrogen bonding.
  \item Interaction fingerprints encoding the residue interaction patterns\cite{CH8_bouysset2021prolif}.
  \item Pose clusters based on ligand heavy-atom RMSD.
  \item Warnings about steric clashes or strained conformations.
\end{itemize}

The main challenge is that none of these quantities is directly observable experimentally. They are proxies for the thermodynamic and kinetic processes that determine drug binding. The main objective of any docking method is for the search algorithm to find biologically reasonable poses. The ability of the scoring function to rank those poses and ligands in a way that correlates with experimental affinity data\cite{CH8_wang2005,CH8_Cheng2009}.

\section{Qualitative Analysis of Docking Poses}

\subsection{Visual Inspection}

Despite extensive automation, visual inspection of docking poses remains indispensable. Most workflows almost always subject top-ranked poses to manual or semi-automated review by an experienced computational biologist before making design or screening decisions\cite{CH8_Kitchen2004,CH8_Charting2015}.

Visual inspection should aim to assess:

\begin{itemize}
  \item \emph{Geometric feasibility}: Does the ligand bind sensibly in the pocket without severe steric clashes or burial of polar atoms in apolar environments without compensation?
  \item \emph{Key pharmacophoric interactions}: Are canonical interactions, such as hinge hydrogen bonds in kinases, catalytic residues in proteases, metal-chelating motifs, satisfied?
  \item \emph{Internal strain}: How are torsion angle preferences in high-energy conformations relative to known low-energy conformers or crystallographic analogues?
  \item \emph{Water molecules}: Are key waters displaced or retained in a manner consistent with known structures and thermodynamic expectations?
\end{itemize}

\subsection{Warning signs to watch for during analysis}

Despite knowing what to expect while analysing docking outcomes, it is seldom discussed what not to accept while visualising docking poses. The following points are non-exhaustive recommendations:

\begin{itemize}
  \item \emph{Buried, unsatisfied charges}: Poses in which a formal charge, such as those from a protonated amine or deprotonated carboxylate group, is completely buried in a hydrophobic pocket without forming salt bridges or strong hydrogen bonds are artefacts, even if highly scored.
  \item \emph{Severe steric clashes}: Ligand atoms overlapping protein atoms, or gross van der Waals clashes between rings and side chains, indicate a failure in the search algorithm or scoring functions. These poses should be discarded, or the protocol re-examined.
  \item \emph{Unrealistic ligand conformations}: Highly twisted ring systems or torsions that deviate dramatically from known crystallographic or QM-optimised conformations, without clear compensating interactions, suggest over-fitting of the scoring function.
  \item \emph{Unrealistic metal coordination}: Poses that ignore established coordination geometries and preferred donor atoms for metal ions, like Zn$^{2+}$ or Mg$^{2+}$, should be treated with great suspicion unless supported by strong experimental evidence.
\end{itemize}

\section{Quantitative Analysis}

Docking scores are often reported in units of binding free energy, such as kcal/mol or kJ/mol; however, they are rarely parameterised to yield quantitatively accurate $\Delta G$ values. Comparative assessments show that simple linear fits between scores and experimental affinities are usually target-specific and often exhibit large residuals\cite{CH8_Warren2006,CH8_Cheng2009}. Therefore:

\begin{itemize}
  \item Treat scores as relative ranking tools within a given target and protocol.
  \item Avoid comparing scores across different targets, protein conformations, or docking algorithms.
  \item Be sceptical of fine distinctions in score, differences $<$ 1-2 kcal/mol between ligands may well be within the error region of the scoring function, and not a true difference in binding affinity.
\end{itemize}

\subsection{Score distributions and rank-based analysis}

Rather than focusing on individual scores, it is more robust to analyse:

\begin{itemize}
  \item \emph{Score distributions}: Histograms of scores for all docked ligands often reveal whether a few ligands are clear outliers or whether the distribution is narrow and undifferentiated.
  \item \emph{Rank-based enrichment:} When known actives are present, one can compute enrichment factors (EF), receiver operating characteristic (ROC) curves, and early enrichment metrics such as BEDROC to evaluate how well the scoring function prioritises actives near the top.
  \item \emph{Score vs physicochemical properties}: Plots of score versus molecular weight, cLogP, polar surface area, or charge can reveal systematic biases, such as those in the scoring functions that favour larger or more hydrophobic compounds.
\end{itemize}

\subsection{What to avoid while reporting the outcome?}

Problematic practices include:

\begin{enumerate}
  \item \emph{Using raw scores as binding free energies}: Statements such as the docking score of -11 kcal/mol indicate nanomolar affinity are not supported by benchmarking data and should be avoided\cite{CH8_Warren2006}.
  \item \emph{Hard score cutoffs without context}: Applying an arbitrary score threshold, such as keep only ligands with score $<$ -8 kcal/mol, without prior validation, can remove true actives or retain large numbers of false positives.
  \item \emph{Cross-target score comparison}: Comparing scores of the same ligand docked to two different proteins, and inferring selectivity directly from the difference is unjustified unless the scoring function has been carefully validated for that purpose.
\end{enumerate}

\section{Internal Validation of Docking Protocols}

\subsection{Re-docking: Reproducing Known Binding Modes}

Re-docking tests use crystallographic complexes with known bound ligands. The ligand is removed, re-docked into the binding site, and the resulting poses are compared to the experimental binding mode, usually via RMSD. A typical acceptance criterion is that at least one of the top-ranked poses has heavy-atom RMSD $<$ 2.0\r{A} from the crystal pose\cite{CH8_Warren2006,CH8_Kitchen2004,CH8_Cheng2009}.

Best practices for re-docking include:

\begin{itemize}
  \item Using multiple co-crystal structures if available, to probe robustness across conformational states,
  \item Evaluating both pose \emph{reproduction} (RMSD) and \emph{ranking} (whether the near-native pose is ranked among the top poses),
  \item Ensuring that protein and ligand preparation protocols used in re-docking match those planned for virtual screening.
\end{itemize}

Only the best RMSD achievable by any pose, while ignoring that the near-native pose was ranked very low; such protocols often fail in screening contexts and must be avoided as they are misleading. 

\subsection{Cross-Docking and Induced Fit Considerations}

Cross-docking, in which ligands are docked into non-cognate protein conformations, assesses sensitivity to receptor flexibility and pocket variability. This is particularly relevant for highly flexible targets such as kinases, nuclear receptors, or GPCRs\cite{CH8_Huang2011,CH8_PandFutureDocking2020}.

Important recommendations:

\begin{itemize}
  \item Include representative open and closed conformations when available.
  \item Quantify success rates across multiple structures, not just a single \emph{best} receptor.
  \item Consider ensemble docking or induced-fit docking when single structure performance is poor.
\end{itemize}

A common mistake is to perform cross-docking across highly different conformational states (e.g., \emph{apo} vs \emph{holo} with large loop movements) without adjusting the protocol or interpreting failures in light of these structural differences.

\subsection{Enrichment-based validation with actives and decoy sets}

When a set of known actives is available, along with an appropriate set of inactive or decoy compounds, virtual screening-style validation can be performed\cite{CH8_Verdonk2004,CH8_DUD-E2012}.

Best practices for enrichment-based validation include:

\begin{itemize}
  \item \emph{Property-matched decoys}: Use decoy sets that match simple physicochemical properties (e.g., molecular weight, cLogP, HBD/HBA counts) but differ in topology from actives; the DUD-E resource is a widely adopted standard\cite{CH8_DUD-E2012}.
  \item \emph{Early recognition metrics}: Compute enrichment factors at low percentages of the ranked list (e.g., EF$_{1\%}$, EF$_{2\%}$), ROC-AUC, and early recognition metrics such as BEDROC\cite{CH8_Truchon2007}.
  \item \emph{Multiple random splits}: When decoys are custom-generated, use multiple random subsets to avoid over-fitting to a single decoy realisation.
\end{itemize}

What not to do includes:

\begin{enumerate}
  \item Using random \emph{drug-like} compounds as decoys without property matching, leading to artificially inflated enrichment,
  \item Reporting only ROC-AUC without considering that high AUC can coexist with poor early enrichment, which is more relevant for experimental follow-up\cite{CH8_Truchon2007},
  \item Validating on the same dataset used for method development, leading to optimistic performance estimates.
\end{enumerate}

\section{Retrospective and prospective benchmarks}

\subsection{Large-scale comparative studies}

Several large-scale comparative studies\cite{CH8_Warren2006,CH8_Cheng2009,CH8_wang2005} have systematically assessed docking programs and scoring functions across diverse targets, highlighting that pose prediction is usually more reliable than affinity ranking. No single docking program or scoring function is universally superior. It was observed that performance varies strongly across target classes and binding-site features Such studies motivate target-specific validation and, in many cases, the use of consensus strategies or re-scoring.

\subsection{Public benchmark sets}

Public benchmarks such as DUD-E, DEKOIS, and others provide standardised actives and decoys across many targets, facilitating method comparison and protocol optimization\cite{CH8_DUD-E2012}. However, users should be aware that benchmarks can still exhibit biases (e.g., analogue bias, decoy selection biases) and may not capture all nuances of a specific industrial target class.

\section{Re-scoring strategies}

\subsection{Consensus scoring}

Consensus scoring combines outputs from multiple docking programs or scoring functions, typically by averaging scores, aggregating ranks, or using more sophisticated schemes such as weighted or exponential rank combinations\cite{CH8_Cheng2009,CH8_ConsensusDocking2022}. Thus, it can reduce method-specific noise and systematic biases. It has been shown to improve enrichment in retrospective virtual screening. It is often straightforward to apply when multiple methods are already in use.
However, it requires careful standardisation of input scores (e.g., sign conventions, scales). A poor averaging or weighting scheme can lead to poor performance or excessive weighting of one scoring method. If additional docking trials are warranted, this increases computational cost.

\subsection{Empirical and knowledge-based Re-scoring functions}

Several standalone scoring functions (e.g., X-Score, DrugScore, PLP variants) were designed explicitly for re-scoring of docked poses, often calibrated on curated protein–ligand complexes\cite{CH8_Cheng2009}. These can be applied to poses generated by any docking engine, decoupling pose generation from scoring.
Recommended practice suggests:

\begin{itemize}
  \item To ensure that the parameterisation domain of the re-scoring function (e.g., target types, ligand chemotypes) is within the domain of application for the intended application.
  \item Validate re-scoring performance using re-docking and enrichment tests on relevant systems.
  \item Examine correlation between original and re-scored ranks to understand how much reordering occurs.
\end{itemize}

\subsection{Physics-based re-scoring}

Physics-based post-processing methods such as MM/PBSA and MM/GBSA approximate binding free energies by combining molecular mechanics energies with continuum solvation models, often using snapshots from molecular dynamics (MD) simulations\cite{CH8_Hou2011,CH8_MMGBSApose2011,CH8_MMGBSAreview2015}. Multiple studies have shown that, when carefully applied, MM/GBSA can outperform standard docking scores in ranking congeneric ligands and in discriminating near-native from decoy poses\cite{CH8_Hou2011,CH8_MMGBSApose2011,CH8_MMGBSAffinity2016,CH8_Xu2013}.

Key considerations for MM/PBSA and MM/GBSA re-scoring include:

\begin{itemize}
  \item \emph{Sampling}: Adequate conformational sampling is essential; short MD trajectories (1–5~ns) can suffice for many systems but may be inadequate for highly flexible binding sites\cite{CH8_Hou2011,CH8_Xu2013}.
  \item \emph{Force field and charge models}: Choice of protein force field and ligand charge model (e.g., RESP vs AM1-BCC) has a substantial impact on predictive performance, and best choices can be system-dependent\cite{CH8_Xu2013,CH8_MMGBSAreview2015}.
  \item \emph{Entropy}: Entropic contributions are often approximated or neglected; relative ranking within congeneric series can still be meaningful even without explicit entropy treatment, however, absolute binding free energies may be less reliable.
\end{itemize}

Practices to avoid while using MM/GBSA or MM-PBSA for re-scoring:

\begin{enumerate}
  \item Reporting MM/GBSA or MM-PBSA values with four decimal points and interpreting them as precise $\Delta G$ values,
  \item Neglecting to specify key parameters (dielectric constants, GB model, number of snapshots, simulation length),
  \item Using MM/GBSA  or MM-PBSA on single, un-minimised docking poses without any relaxation, which can amplify artefacts from docking.
\end{enumerate}

\subsection{MD-based refinement}

MD simulations can be used to relax docking poses and explore local conformational space, followed by checking for hydrogen bonds, hydrophobic contacts, and water networks that are only crudely approximated in docking\cite{CH8_meng2011,CH8_MMGBSApose2011}.

Practical suggestions:

\begin{itemize}
  \item Focus MD refinement on a manageable subset of promising ligands (e.g., tens to low hundreds),
  \item Monitor RMSD of ligand and nearby residues to ensure that the complex remains physically reasonable,
  \item Analyse occupancy of key waters, especially in cases where \emph{displaceable} vs  \emph{conserved} waters may influence potency.
\end{itemize}

\section{Best-Practice Workflow}

\subsection{Protein and Ligand Preparation}

Robust protein and ligand preparation is foundational. Recommended steps include:

\begin{itemize}
  \item Choose high-quality experimental structures, preferably in complex with ligands similar to the intended chemotypes;
  \item Assign protonation states and tautomers for ionisable residues and ligands at relevant pH using established tools; consider alternative states for borderline cases;
  \item Model missing loops and side chains where necessary, however, avoid over-modelling when structural uncertainty is high.
  \item Retain or remove crystallographic waters based on well-defined criteria (e.g., buried, conserved waters vs bulk-like, weakly bound waters).
\end{itemize}

Not so good practices that have been observed are:

\begin{enumerate}
  \item Using an \emph{apo} structure with a collapsed binding pocket for docking without any attempt to open or relax it.
  \item Ignoring biologically relevant cofactors or metal ions (e.g., Mg$^{2+}$ in kinases, Zn$^{2+}$ in metalloproteases),
  \item Assuming all titratable groups are in their default protonation state at pH 7, regardless of micro-environment.
\end{enumerate}

\subsection{Docking protocol design}

A typical docking campaign might proceed as follows:

\begin{enumerate}
  \item \emph{Protocol validation}: Perform re-docking and enrichment-based tests on representative co-crystal ligands and in-house actives, adjusting parameters (grid size, sampling exhaustiveness, constraints) to achieve acceptable pose recovery and enrichment.
  \item \emph{Primary scoring}: Dock a large library with high-throughput settings, using a validated scoring function and retaining multiple poses per ligand.
  \item \emph{Post-processing}: Apply property and substructure filters (e.g., removal of PAINS and undesirable reactive groups), cluster by chemotype, and perform qualitative pose inspection of top-ranked chemotypes.
  \item \emph{Re-scoring}: Apply consensus scoring and/or MM/GBSA re-scoring to a reduced set, followed by deeper visual curation.
  \item \emph{Selection for synthesis or purchase}: Prioritise compounds based on a combination of score, pose quality, novelty, synthetic tractability, and alignment with project hypotheses.
\end{enumerate}

\section{Common Pitfalls}

\subsection{Score over-Interpretation}

One of the most pervasive errors is to treat docking scores as if they were experimentally measured affinities. Examples include:

\begin{itemize}
  \item Ranking analogues solely based on small score differences and proposing SAR trends without any validation.
  \item Inferring potency categories (e.g., micromolar vs nanomolar) directly from score ranges.
  \item Claiming that a ligand \emph{binds more strongly} than a reference compound purely because its docking score is more negative by 0.5-1.0 units.
\end{itemize}

Such claims are not supported by benchmarking data and often contradict observed uncertainties in scoring functions\cite{CH8_Warren2006,CH8_Cheng2009}.

\subsection{Ignoring protonation, tautomerism, and stereochemistry}

Another frequent failure mode is to dock a single arbitrary protonation state, tautomer, or stereoisomer of a ligand and then over-interpret the resulting pose and score. Bad practices include:

\begin{enumerate}
  \item Docking a ligand with multiple basic centres in a fully protonated form in a hydrophobic pocket,
  \item Ignoring alternative tautomers that would better complement the hydrogen bond pattern of the pocket,
  \item Neglecting to consider all stereoisomers when chirality is not experimentally defined.
\end{enumerate}

Best practice is to enumerate chemically reasonable protonation states, tautomers, and stereoisomers, filter them based on pH and prior knowledge, and dock the relevant set, assessing whether the binding mode is robust across plausible states.

\subsection{Cherry-picking poses and ligands}

Docking generates many poses; it is tempting to select only those that fit a desired hypothesis or structural narrative. Problematic behaviours include:

\begin{itemize}
  \item Highlighting a pose that matches a pre-conceived pharmacophore while ignoring that it is ranked very low or exhibits clashes,
  \item Presenting only a few highly favourable scores while omitting the broader distribution and the presence of close analogues with poor scores,
  \item Modifying the protein structure or protonation states post hoc to accommodate a particular ligand without systematically re-validating the protocol.
\end{itemize}

These practices bias interpretation and undermine the predictive value of the calculations.

\subsection{Treating docking as an absolute proof of binding}

Docking is a hypothesis-generation tool, not proof of binding or mechanism. Observed bad practices include:

\begin{enumerate}
  \item Asserting that docking \emph{confirms} experimental binding or mechanism solely because a plausible pose can be found,
  \item Using docking alone to claim that a compound is selective for one target over another,
  \item Drawing mechanistic conclusions (e.g., about covalent bond formation or proton transfer) from docked structures that do not include the necessary quantum-chemical treatment.
\end{enumerate}

A more appropriate stance is to treat docking as \emph{supporting} evidence, to be tested alongside experimental data and other computational methods.

\subsection{Misuse of decoy dets and benchmarks}

Finally, validation itself can be misused. Examples include:

\begin{itemize}
  \item Tuning a scoring function or protocol on a benchmark set (e.g., DUD-E) and then claiming superior performance based on the same set, without external validation\cite{CH8_DUD-E2012}.
  \item Constructing decoy sets that are trivially separable from actives (e.g., by size or polarity), thereby inflating enrichment metrics,
  \item Comparing methods using different decoy sets or evaluation metrics, leading to incommensurate results.
\end{itemize}

Best practice is to use independent training and test sets, employ property-matched decoys, and report multiple, complementary performance metrics.

\section{Summary}

In a nutshell, the chapter presents a critical framework for interpreting and validating molecular docking results, highlighting that the analysis stage represents the weakest link in docking workflows. While docking methods can often reproduce experimental binding poses, their capacity to rank ligands by binding affinity remains limited and is highly dependent on the specific system. Qualitative inspection is crucial for reliable interpretation of docking results. Valid docking poses should demonstrate geometric correctness, appropriate pharmacophoric interactions, minimal strain, and be consistent with established structural features. From a quantitative perspective, docking scores are most appropriately used for relative ranking within a single, validated protocol. Minor differences in scores are typically insignificant, and comparisons across different targets are unreliable. Key validation strategies include re-docking, cross-docking, and enrichment studies using property-matched decoys. Emphasis is placed on the necessity of target-specific validation and the need for caution to avoid overfitting and improper use of benchmark datasets. Post-processing strategies, such as consensus scoring, MM/GBSA or MM-PBSA, and molecular dynamics-based refinement, can improve ranking accuracy and structural reliability. These methods require careful implementation and sufficient sampling to ensure validity, and they primarily serve as tools for generating hypotheses. The validity of the conclusions depends on integration with experimental data, thorough validation, and avoiding common pitfalls, such as over-interpretation of scores, neglecting chemical intuitions, and selective reporting of favourable results.

\bibliographystyle{unsrtnat}
\bibliography{chapter8}

\end{document}